% ==============================================================================
% Chapter 1: Introduction & Problem Formulation
% Target Page Budget: ~5 pages
% ==============================================================================
\chapter{Introduction \& Problem Formulation}
\label{chap:introduction}

\section{Background and Societal Context}
\label{sec:intro_background}
% [Drafting Guidance: ~1.5 pages]
% - Pervasiveness of digital media platforms (Instagram, TikTok, YouTube).
% - Transition from organic social networking to algorithmic recommendation engines.
% - The hyper-curation of lifestyle imagery and its cognitive footprint.
% - Initial citations: \cite{twenge2018increases, primack2017social}.

\section{The Mechanics of Digital Comparison and Influencer Culture}
\label{sec:intro_influencer_comparison}
% [Drafting Guidance: ~1.0 page]
% - The rise of commercial influencer marketing and parasocial relationship formation.
% - Blurring boundaries between authentic peer sharing and sponsored commercial messaging \cite{lou2019influencer, deveirman2017marketing}.
% - Unconscious upward social comparison triggers in algorithmic feeds.

\section{Problem Statement and Core Research Challenges}
\label{sec:intro_problem_statement}
% [Drafting Guidance: ~1.0 page]
% - Core challenge: Lack of cognitive transparency in how feeds impact individual well-being in real-time.
% - Passive scrolling vs. active engagement dilemma.
% - The limitations of existing digital detox tools (crude screen-time limits vs. cognitive awareness).

\section{Research Questions and Guiding Hypotheses}
\label{sec:intro_research_questions}
% [Drafting Guidance: ~0.75 pages]
% Formulate primary research questions (RQ1-RQ4):
% RQ1: How can multi-agent systems reliably detect and quantify upward social comparison triggers in multimodal social media streams?
% RQ2: What cognitive state modeling techniques best capture longitudinal shifts in user well-being?
% RQ3: Can autonomous reflection agents provide timely, non-intrusive cognitive scaffolding without triggering reactance?
% RQ4: What ethical boundaries must govern autonomous affective monitoring agents?

\section{Scope, Project Objectives, and Anticipated Contributions}
\label{sec:intro_contributions}
% [Drafting Guidance: ~0.5 pages]
% - Delineate the scope for Assignment 2 (Preparatory System Specification).
% - List of tangible contributions: Theoretical framework, Agent society specification, Simulation protocol, Ethical roadmap.

\section{Structure of the Preparatory Report}
\label{sec:intro_outline}
% [Drafting Guidance: ~0.25 pages]
% - Brief chapter-by-chapter roadmap guiding the reader through Chapters 2 to 6.
This report is organized into six primary chapters:
\Cref{chap:literature_review} synthesizes foundational psychology and AI agent literature;
\Cref{chap:system_architecture} details the system architecture and multi-agent roles of \projectname{};
\Cref{chap:methodology} presents the experimental simulation framework and metrics;
\Cref{chap:implementation_roadmap} outlines technical feasibility and implementation milestones; and
\Cref{chap:conclusion} concludes with future perspectives and impact.
